% 配置文件：填写参赛项目信息
% ============================================

% ---------- 参赛基本信息 ----------
\newcommand{\competitionYear}{2025}           % 竞赛年份
\newcommand{\competitionSession}{十八}        % 第几届
\newcommand{\projectTitle}{项目名称}           % 题目名称
\newcommand{\projectSubType}{}               % 题目子类型（仅题目5、6、9需填写）

% ---------- 参赛题目选择 ----------
% 已选中的题目用 \checkmark 表示，未选中留空 {}
% 用法：\newcommand{\topicXxx}{\checkmark} 或 \newcommand{\topicXxx}{}

% 命题类
\newcommand{\topicOne}{}          % 题目1：微小位移测量
\newcommand{\topicTwo}{}          % 题目2：探究电磁感应现象中的能量转换
\newcommand{\topicThree}{}        % 题目3：弱压力测量
\newcommand{\topicFour}{}         % 题目4：晶体双折射

% 制作类
\newcommand{\topicFive}{}         % 题目5：实验仪器的制作、改进
\newcommand{\topicFiveSubA}{}     % 题目5子类A：定量类
\newcommand{\topicFiveSubB}{}     % 题目5子类B：演示类
\newcommand{\topicSix}{}          % 题目6：前沿物理
\newcommand{\topicSixSubA}{}      % 题目6子类A：前沿文献类
\newcommand{\topicSixSubB}{}      % 题目6子类B：自研文献类
\newcommand{\topicSeven}{\checkmark}  % 题目7：AI+物理实验

% 信息资源类
\newcommand{\topicEight}{}        % 题目8：大学物理教学微视频
\newcommand{\topicNine}{}         % 题目9：物理教学资源开发
\newcommand{\topicNineSubA}{}     % 题目9子类A：媒体资源类
\newcommand{\topicNineSubB}{}     % 题目9子类B：仿真程序类

% 注：\checkmark 定义为 \ding{51}（✓），需要 pifont 宏包
%     如需手动在 □ 中画 √，请使用原始 □ 符号（U+25A1）配合手写标记
